\documentclass[10pt,a4paper]{amsart}
\usepackage[margin=19mm]{geometry}
\usepackage{amsmath,amssymb,mathtools}
\usepackage[scr=rsfs,cal=euler]{mathalfa}
\usepackage{xfrac}
\usepackage[unicode,hidelinks,bookmarks=false]{hyperref}
\newcommand{\ie}{i.e.\ }
\newcommand{\cf}{cf.\ }
\newcommand{\resp}{respectively\ }
\newcommand{\Subsection}{\subsection}
\providecommand{\nobreakdash}{\nobreak\hbox{-}}
\setlength{\emergencystretch}{2em}
\multlinegap0pt
\usepackage[all]{xy}
\usepackage{amssymb}
\newtheorem{proposition}{Proposition}
\newcommand{\EM}{\mathbf{EM}}
\newcommand{\PEM}{\mathbf{PEM}}
\newcommand{\bbT}{\mathbb{T}}
\title{Projective error models: Stabilizer codes, Clifford codes, and weak stabilizer codes}
\author{Jonas Eidesen}
\date{Source: arXiv:2506.01843v3 (CC BY 4.0)}
\begin{document}
\maketitle

\noindent\emph{Source and licence.} Projective error models: Stabilizer codes, Clifford codes, and weak stabilizer codes. Version arXiv:2506.01843v3 (CC BY 4.0), distributed by arXiv under the Creative Commons Attribution 4.0 International licence, http://creativecommons.org/licenses/by/4.0/, as recorded in the arXiv licence metadata for this exact version and shown on the abstract page https://arxiv.org/abs/2506.01843v3. Full licence notice: \texttt{licenses/arxiv-2506.01843-CC-BY-4.0.txt}.\par\smallskip

\noindent\textit{Editorial scope.} Counted unit: the article's proposition PEM (source0.tex lines 345--349). The article's complete original proof of it is delivered with it (source0.tex lines 350--399, one \texttt{proof} environment of 50 lines). No mathematics is written, repaired or re-derived: the statement and the proof are the article's own lines, in the article's own notation, and no retained line is substituted or re-flowed.  No further source-local context is needed: every cross-reference of every retained line resolves inside the two delivered spans.  Nothing was omitted from any retained span, so the package carries no omission record.  No retained line contains a citation, so the wrapper carries no bibliography and no bibliography entry.  No retained line contains a figure environment, an image-inclusion command or drawing code, so no image is delivered and the package declares no asset dependency.  Editorial formatting only, in addition: an AMS \texttt{amsart} preamble that restates the article's own declarations and macros used by the retained lines, namely the environments \texttt{proposition} on separate counters, together with the standard packages those lines need.  The retained environments print in this document as Proposition 1 in order of appearance, whereas the article prints the same items with the numbering of its own full document.  The complete native source of the article as distributed in the arXiv e-print for this exact version is bundled unchanged as \texttt{sources/179175/source0.tex}.
\begin{proposition}\label{prop:correspondence between EM and PEM}
    Let $V$ be a Hilbert space. If $(E,\lambda) \in \EM_V$ then the quotient $G := E/Z(E)$ admits a projectively faithful irreducible projective representation $\pi$ on $V$. Furthermore, the group homomorphism $q \circ \pi$ is uniquely determined by the representation $\lambda$.

    Conversely, if $(G',\pi') \in \PEM_V$, then for any natural number $n$ such that the order of $[\sigma_{\pi'}] \in H^2(G_2,\bbT)$ divides $n$, we have that there exists a central extension $E'$ of $G'$ by $C_n$ that admits a faithful irreducible representation on $V$.
\end{proposition}

\begin{proof}
    For the first statement, let $(E,\lambda) \in \EM_V$. We will define a projective representation of $G$ on $V$. For each $x \in G$ let $s_x \in E$ be any coset representative of $x$. For $x,y \in G$ we have that there exists a unique element $\alpha(x,y) \in Z(E)$ such that $s_x s_y = \alpha(x,y) s_{xy}$. Since $\lambda$ is an irreducible representation of $E$ we have by Schur's Lemma that
    \begin{equation*}
        \lambda(\alpha(x,y)) = \sigma(x,y)1_{U(V)}
    \end{equation*}
    for some $\sigma(x,y) \in \bbT$. Letting $x$ and $y$ range over all elements in $G$ we obtain a 2-cocycle $\sigma \in Z^2(G,\bbT)$. We now define a function $\pi \colon G \to U(V)$ by the equation
    \begin{equation*}
        \pi(x) := \lambda(s_x).
    \end{equation*}
    Then $\pi$ is a $\sigma$-projective representation of $G$. Note that by construction we have that the diagram
    \begin{equation*}
        \xymatrix{
            E \ar[r]^-{p} \ar[d]_-{\lambda_1} & G \ar[d]^-{q \circ \pi_1} \\
            U(V) \ar[r]^-{q} & PU(V)
        }
    \end{equation*}
    commutes. By surjectivity of the quotient map $p \colon E \to G$ we get that $q \circ \pi$ is the unique group homomorphism making this diagram commute.
    
    It remains to check that $\pi$ is projectively faithful: Let $x \in \ker{(q \circ \pi)}$. By assumption we have that 
    \begin{equation*}
        \lambda(s_x) = \pi(x) = z 1_{U(V)}
    \end{equation*}
    for some $z \in \bbT$. Hence, for any $g \in E$ we have that
    \begin{equation*}
        \lambda(s_x g) = \lambda(g s_x).
    \end{equation*}
    Since $\lambda$ is injective ($\lambda$ is a faithful representation of $E$ on $V$) we have that $s_x g = g s_x$ for every $g \in E$. Thus, $s_x \in Z(E)$, meaning that $x = 1_{G}$. Hence, $q \circ \pi$ is injective.

    For the second statement, let now $(G',\pi') \in \PEM_V$, and suppose that $\pi'$ is a $\sigma$-projective representation of $G'$ on $V$. Let $n$ be a natural number such that the order of $[\sigma] \in H^2(G',\bbT)$ divides $n$. The following short exact sequence,
    \begin{equation*}
        \xymatrix{
            1 \ar[r] & C_n \ar[r]^-{i} & \bbT \ar[r]^-{(-)^n} & \bbT \ar[r] & 1,
        }
    \end{equation*}
    induces a long exact sequence in group cohomology with the following segment,
    \begin{equation*}
        \xymatrix{
            \cdots \ar[r] & H^2(G',C_n) \ar[r]^-{i^*} & H^2(G',\bbT) \ar[r]^-{n} & H^2(G',\bbT) \ar[r] & \cdots,
        }
    \end{equation*}
    where $n$ here denotes the map that is multiplication by $n$. The fact that the order of $[\sigma]$ divides $n$ is precisely the fact that $[\sigma] \in \ker{n}$, and since this sequence is exact, this means that there exists an element $[\sigma'] \in H^2(G',C_n)$ such that $i^*([\sigma']) = [\sigma]$. Fix a representative $\sigma' \colon G' \times G' \to C_n$ of $[\sigma']$ and choose a function $f \colon G' \to \bbT$ such that $\sigma' = (\delta f) \, \sigma$. Then we have that
    \begin{equation*}
        f(x)\pi'(x) \, f(y)\pi'(y) = \sigma'(x,y) f(xy)\pi'(xy) \text{ for all } x,y \in G'.
    \end{equation*}
    Let $E'$ denote the central extension of $G'$ by $C_n$ corresponding to the element $[\sigma'] \in H^2(G',C_n)$. Then define a representation $\lambda' \colon E' \to U(V)$ by
    \begin{equation*}
        \lambda'(z,x) = z f(x)\pi'(x)
    \end{equation*}
    for $z \in C_n$ and $x \in G'$. This is easily seen to define a faithful irreducible representation of $E'$ on $V$, completing the proof.
\end{proof}

\end{document}
